% TOP_PROBLEM: {"id": 1896, "title": "Erdős Problem #39", "category_id": 2, "category": {"id": 2, "name": "combinatorics", "display_name": "Combinatorics", "description": "Counting problems, graph theory, discrete structures.", "slug": "combinatorics", "order_index": 2, "created_at": "2026-07-31T15:26:25.671Z"}, "status": "open"}
% FIRST_POSTED: null
% ENTRY_KIND: statement_only
% Imported from: unsolvedmath_additional_500.tex; UnsolvedMath ID 1896
% !TeX program = lualatex
% Compile twice: lualatex 1896.tex
% Self-contained: no external bibliography, images, or \input files.
% Numeric section labels and visible numbers are original UnsolvedMath IDs.
\documentclass[11pt,a4paper]{article}
\usepackage[margin=24mm,headheight=14pt]{geometry}
\usepackage{fontspec}
\setmainfont{Latin Modern Roman}
\setsansfont{Latin Modern Sans}
\setmonofont{Latin Modern Mono}
\usepackage{amsmath,amssymb,mathtools}
\usepackage{unicode-math}
\setmathfont{Latin Modern Math}
\usepackage{microtype}
\usepackage{enumitem,xurl,needspace}
\usepackage[dvipsnames]{xcolor}
\usepackage{hyperref}
\hypersetup{unicode=true,colorlinks=true,linkcolor=MidnightBlue,urlcolor=MidnightBlue,
 pdftitle={UnsolvedMath Problem 1896},pdfauthor={Source-annotated compilation},bookmarksnumbered=true}
\usepackage{bookmark,tocloft,titlesec,fancyhdr}
\setlength{\cftsecnumwidth}{4.2em}
\cftsetpnumwidth{2.5em}
\cftsetrmarg{4em}
\titleformat{\section}{\Large\bfseries\raggedright}{\thesection}{0.7em}{}
\pagestyle{fancy}\fancyhf{}
\fancyhead[L]{\small\sffamily Additional UnsolvedMath statements}
\fancyhead[R]{\small Original catalogue IDs}
\fancyfoot[C]{\thepage}
\setlength{\parindent}{0pt}
\setlength{\parskip}{5pt plus 1pt}
\setlength{\emergencystretch}{3em}
\setcounter{tocdepth}{1}\setcounter{secnumdepth}{1}
\setlist[itemize]{leftmargin=3em,itemsep=4pt,topsep=4pt,parsep=0pt}
\allowdisplaybreaks
\newcommand{\N}{\mathbb{N}}
\newcommand{\Z}{\mathbb{Z}}
\newcommand{\Q}{\mathbb{Q}}
\newcommand{\R}{\mathbb{R}}
\newcommand{\C}{\mathbb{C}}
\newcommand{\F}{\mathbb{F}}
\newcommand{\A}{\mathbb{A}}
\newcommand{\PP}{\mathbb{P}}
\newcommand{\E}{\mathbb{E}}
\newcommand{\eps}{\varepsilon}
\newcommand{\dd}{\,\mathrm d}
\newcommand{\sref}[2]{\hyperlink{source:#1:#2}{[S#2]}}
\newif\ifshowcatalogue
\showcataloguetrue
\newcommand{\cataloguescope}[1]{\ifshowcatalogue
 \paragraph{Statement recorded in the supplied source.}
 #1\fi}
\begin{document}
% UnsolvedMath ID 1896; source catalogue code EP-39

\setcounter{section}{1895}

\section{An Infinite Sidon Set of Near Square-Root Growth}\label{1896}

{\small\sffamily UnsolvedMath ID 1896 · Combinatorics}

\subsection{Definitions and mathematical statement}

\paragraph{Shared notation.}Unless a section says otherwise, $\N=\{1,2,\ldots\}$,
$\Z$ is the ring of integers, $\Q,\R,\C$ are the rational, real, and complex
fields, $[N]=\{1,\ldots,\lfloor N\rfloor\}$, and $\log$ is the natural
logarithm. For real functions of a parameter tending to infinity,
$f\ll g$ and $f=O(g)$ mean $|f|\leq Cg$ eventually for some fixed $C$;
$f\gg g$ reverses this comparison; $f=o(g)$ means $f/g\to0$;
$f\sim g$ means $f/g\to1$; and $f\asymp g$ means both order bounds.
A subscript on $O$ or $\ll$ specifies allowed constant dependence.
The expression $n^{o(1)}$ in an upper bound means growth smaller than
$n^\epsilon$ for each fixed $\epsilon>0$. Local definitions govern any
exception, finite-field convention, or use of the same symbol for another
object. An asymptotic claim is not automatically an effective algorithm.



A Sidon set has $a+b=c+d$ only when the multisets $\{a,b\}$ and $\{c,d\}$ agree. Equal summands are allowed, but interchanging their order is not a new representation. Unless a real or group domain is explicitly specified, the set here consists of integers. \sref{1896}{1} \sref{1896}{2}

\paragraph{Statement recorded in the supplied source.}
Is there an infinite Sidon set $A\subset \mathbb{N}$ such that $ \lvert A\cap \{1\ldots,N\}\rvert \gg_\epsilon N^{1/2-\epsilon} $ for all $\epsilon>0$? \sref{1896}{1}

\subsection{Short English statement}

Can one infinite Sidon set remain denser than every fixed power below the square-root scale throughout all sufficiently large intervals?

\subsection{Sources}

{\small

\begin{itemize}

\item[\hypertarget{source:1896:1}{S1}] ULAM.AI, \emph{UnsolvedMath}, supplied \texttt{problems.json}, record 1896 (EP-39), statement and background. The embedded literature-review date is 2026-08-17; that is the archive’s review date, not a fresh certification. Dataset landing page: \url{https://huggingface.co/datasets/ulamai/UnsolvedMath}.

\item[\hypertarget{source:1896:2}{S2}] T. F. Bloom, \emph{Erdős Problems}, Problem 39, \url{https://www.erdosproblems.com/39}. Reference imported from the supplied record; the live problem page was not independently checked for this compilation.

\item[\hypertarget{source:1896:3}{S3}] Imre Z. Ruzsa, An infinite Sidon sequence, Journal of Number Theory 68 (1998), 63-71. \url{https://doi.org/10.1006/jnth.1997.2192}. Bibliographic information imported from the supplied record.

\item[\hypertarget{source:1896:4}{S4}] Javier Cilleruelo, Infinite Sidon sequences, Advances in Mathematics 255 (2014), 474-486. \url{https://doi.org/10.1016/j.aim.2014.01.011}. Bibliographic information imported from the supplied record.

\item[\hypertarget{source:1896:5}{S5}] Ajtai, Mikl\'os and Koml\'os, J\'anos and Szemer\'{e}di, Endre, A dense infinite {S}idon sequence. European J. Combin. (1981), 1--11. Bibliographic information imported from the supplied record.

\item[\hypertarget{source:1896:6}{S6}] Erd\H{o}s, P., Problems and results on combinatorial number theory. A survey of combinatorial theory (Proc. Internat. Sympos., Colorado State Univ., Fort Collins, Colo., 1971) (1973), 117-138. Bibliographic information imported from the supplied record.

\end{itemize}

}

\label{end:1896}
\end{document}
